\documentclass[10pt,a4paper]{amsart}
\usepackage[margin=19mm]{geometry}
\usepackage{amsmath,amssymb,mathtools}
\usepackage[scr=rsfs,cal=euler]{mathalfa}
\usepackage{xfrac}
\usepackage[unicode,hidelinks,bookmarks=false]{hyperref}
\newcommand{\ie}{i.e.\ }
\newcommand{\cf}{cf.\ }
\newcommand{\resp}{respectively\ }
\newcommand{\Subsection}{\subsection}
\providecommand{\nobreakdash}{\nobreak\hbox{-}}
\setlength{\emergencystretch}{2em}
\multlinegap0pt
\usepackage{bm}
\usepackage{bbm}
\usepackage{dsfont}
\usepackage{xspace}
\usepackage{cancel}
\usepackage{mathtools}
\usepackage{cleveref}
\usepackage{amsthm}
\makeatletter
\@ifundefined{theorem}{\newtheorem{theorem}{Theorem}}{}
\@ifundefined{definition}{\newtheorem{definition}[theorem]{Definition}}{}
\@ifundefined{lemma}{\newtheorem{lemma}[theorem]{Lemma}}{}
\makeatother
\providecommand{\tightlist}{\setlength{\itemsep}{0pt}\setlength{\parskip}{0pt}}
\title[New lower bounds for the degree/diameter problem via interac]{New lower bounds for the degree/diameter problem via interaction with a browser-accessible LLM}
\author{Ryosuke Mizuno}
\date{Source: arXiv:2606.15860v1 (CC BY 4.0)}
\begin{document}
\maketitle

\noindent\emph{Source and licence.} New lower bounds for the degree/diameter problem via interaction with a browser-accessible LLM. Version arXiv:2606.15860v1 (CC BY 4.0), distributed under the Creative Commons Attribution 4.0 International licence, as recorded in the arXiv licence metadata for this exact version and shown on the abstract page https://arxiv.org/abs/2606.15860v1. Full rights notice: licenses/arxiv-2606.15860-CC-BY-4.0.txt.\par\smallskip

\noindent\textit{Editorial scope.} Counted unit: the Lemma whose source label is \texttt{lem:lift} in the article (source0.tex lines 1387--1404, source label \texttt{lem:lift}), delivered together with the article's own complete original proof of it (source0.tex lines 1406--1449, one \texttt{proof} environment of 44 lines). No mathematics is written, repaired or re-derived: statement and proof are the article's own text line for line. Required context retained uncounted: def:lift (lines 1347--1361). Every definition, display and label of those lines is retained unchanged. Omitted from this package and not redistributed: 1--1346, 1362--1386, 1405--1405, 1450--1899. Nothing inside a retained excerpt is omitted or altered: every retained body is delivered exactly as the article's lines are written, so no omission receipt of any kind accompanies this package. Citations: the retained excerpts carry no citation command at all, so no bibliography entry of the article is transcribed and no reference is redistributed. Reference adaptation: none was needed and none was made; every reference inside the delivered unit resolves inside it, so no unresolved reference or citation is produced. Printed slips kept literal: no printed word, symbol, hypothesis or number of the counted statement or of its proof was altered. Layout and class: the article's own class is \texttt{article} with packages including geometry, amsmath, amssymb, amsthm, array, booktabs, longtable, microtype, natbib, hyperref, cleveref, fontenc, lmodern; the wrapper is the lane's pinned \texttt{amsart} editorial wrapper at 10pt on A4, which restates the theorem-like environments the retained text needs and adds the article's own macro definitions that the retained text needs (\texttt{\textbackslash tightlist}), with extra wrapper packages (bm, bbm, dsfont, xspace, cancel, mathtools, cleveref). The delivered numbering is the unit's own. Assets: no retained span contains a figure environment, a graphics-inclusion command, a PDF-page inclusion, a table or a TikZ picture; no image file is copied into this package, the package declares no asset dependency and no image file is redistributed with it. Licence: version arXiv:2606.15860v1 of this article is distributed under the Creative Commons Attribution 4.0 International licence, as recorded in the arXiv licence metadata for this exact version and shown on the abstract page https://arxiv.org/abs/2606.15860v1. A full rights notice is delivered at licenses/arxiv-2606.15860-CC-BY-4.0.txt. Characters: every retained excerpt is pure ASCII, so the delivered engine, XeLaTeX with its default Latin Modern text font, reports no missing character.
\begin{definition}\label{def:lift}
Let \(C\) be a \(\rho\)-regular simple
undirected graph, \(S\) a symmetric alphabet with \(|S|=\rho\),
\(\ell:\overrightarrow E(C)\to S\) an inverse-consistent labeling, and
\((A_t,b_t)_{t\in S}\) a route chart on \(\mathbb F_q^s\). Define the
route-chart lift graph \[
  G=G(C,S,\ell,A,b)
\] as follows. The vertex set is \[
  V(G)=V(C)\times\mathbb F_q^s .
\] For an undirected edge \(\{u,v\}\) of \(C\) with \(\ell(u\to v)=t\),
for each \(x\in\mathbb F_q^s\) and each \(\lambda\in\mathbb F_q\), add
the undirected edge joining the two vertices \[
  (u,x),\qquad (v,A_tx+\lambda b_t) .
\]
\end{definition}

\begin{lemma}\label{lem:lift}
With the above notation, suppose the route chart is
universal. Then the following hold.

\begin{enumerate}
\def\labelenumi{\arabic{enumi}.}
\tightlist
\item
  \(G\) is a simple undirected graph.
\item
  \(G\) is \(\rho q\)-regular.
\item
  If \(C\) is an exact-NB-\(s\) controller then \(G\) is connected and
  \(\operatorname{diam}(G)\le s\).
\item
  \(|V(G)|=|V(C)|q^s\).
\end{enumerate}
\end{lemma}

\begin{proof}
The graph is undirected by \cref{def:lift}. We first
prove simplicity and the degree. Since \(C\) is simple, for a directed
edge \(u\to v\) of \(C\) we always have \(u\ne v\). Hence \(G\) has no
loops. By universality, each \(b_t\) is nonzero. Indeed, the word
\(t^s\) is reduced, and the last column of the corresponding \(M_{t^s}\)
is \(b_t\); if \(M_{t^s}\) is nonsingular then \(b_t\ne0\).

Consider a fixed vertex \((u,x)\in V(G)\). There are \(\rho\) directed
edges leaving \(u\), and for each, choosing \(\lambda\in\mathbb F_q\)
yields the candidate neighbor \[
  (v,A_tx+\lambda b_t) .
\] By the argument at the end of \cref{def:lift}, these are all the
neighbors of \((u,x)\in V(G)\). Distinct directed edges go to distinct
controller vertices \(v\) because \(C\) is simple. For the same directed
edge, if \[
  A_tx+\lambda b_t=A_tx+\lambda' b_t ,
\] then \((\lambda-\lambda')b_t=0\), and from \(b_t\ne0\) we get
\(\lambda=\lambda'\). Hence \((u,x)\) has exactly \(\rho q\) distinct
neighbors. Therefore \(G\) is simple and \(\rho q\)-regular.

Next we show the diameter. Take any \((c,x),(d,y)\in V(G)\). Since \(C\)
is exact-NB-\(s\), there is an NB walk of length \(s\) in \(C\), \[
  c=c_0,c_1,\ldots,c_s=d .
\] Let the label of \(c_{i-1}\to c_i\) be \(t_i\). As shown in the
previous subsection, under the inverse-consistent labeling, the label
sequence \(w=t_1\cdots t_s\) of an NB walk is a reduced word.

Choose control parameters \(\lambda_1,\ldots,\lambda_s\in\mathbb F_q\)
along this walk. Starting the fiber coordinate at \(x_0=x\) and setting
\[
  x_i=A_{t_i}x_{i-1}+\lambda_i b_{t_i}
  \qquad(1\le i\le s) ,
\] we obtain inductively \[
  x_s=A_wx+M_w
  \begin{pmatrix}\lambda_1\\ \vdots\\ \lambda_s\end{pmatrix} .
\] By universality \(M_w\) is nonsingular, so the linear equation \[
  M_w
  \begin{pmatrix}\lambda_1\\ \vdots\\ \lambda_s\end{pmatrix}
  =y-A_wx
\] has a solution. Using this solution, we obtain a walk of length \(s\)
from \((c,x)\) to \((d,y)\). In particular \(G\) is connected, and \(\operatorname{diam}(G)\le s\). The
vertex count follows immediately from the definition.
\end{proof}

\end{document}
